\begin{afterword}
\summaryhead

The post guesses, and says twice that it is only a guess, that Keats did not understand the explanation of the rainbow he complained about in \textsc{Lamia}. Perhaps he had only read that Newton explained it as light reflected from raindrops, which was known centuries before; Newton's addition was that white light is made of colours. The evidence offered is a toast: Keats, with Lamb, Wordsworth and Haydon, drank ``Confusion to the memory of Newton'' because Newton ``destroyed the poetry of the rainbow by reducing it to a prism.''

If so, Keats had a fake reduction: he had been told that the rainbow was reduced, but had not seen it. The author says anti-reductionists in general confuse professing that something is reducible with seeing it, and that this probably accounts for much of the emptiness said to come with reductionism. Being told that science explained the rainbow only moves it into a genre labelled boring. The post ends with a counter-toast: ``To the memory of Keats's confusion.''

\responsehead

The distinction the post draws is real and useful. Being told that a thing has been explained is not the same as seeing the explanation, and the post connects this to the author's earlier posts on knowledge that is only recited. Its account of Keats's line also fits the poem well: the rainbow ``given / In the dull catalogue of common things'' because ``We know her woof'' is a rainbow moved into a boring genre by report. The history of the rainbow's explanation is, in substance, right.

The case against Keats is weak, and the post marks most of it as guesswork. It depends on two guesses, each marked as a guess, and one piece of evidence, the toast. The toast's wording comes from a letter Haydon wrote to Wordsworth in 1842, nearly twenty-five years after the dinner. Haydon's autobiography tells the story differently: Lamb, ``exceedingly merry,'' mocks Newton, Keats agrees about the rainbow, and everyone drinks ``Newton's health, and confusion to mathematics.'' In both versions the toast is a dinner-table joke. It shows at most an attitude to Newton's rainbow, not what Keats knew about optics. The post offers it only as ``one reason to suspect,'' and it is a weak one.

The step from Keats to anti-reductionists in general has less support still. The claim that fake reduction ``probably accounts for a hell of a lot'' of the emptiness said to come with reductionism is hedged, but the only case offered is the guess about Keats. The closing statement, ``That's how anti-reductionists experience `reductionism','' quotes no anti-reductionist and no literary critic describing that experience. As in the previous post, the opponents are not identified.

In short: a real distinction between being told and seeing, attributed to Keats on the basis of two guesses and a dinner-party toast recalled nearly twenty-five years later, and extended to anti-reductionists in general with no one quoted.
\end{afterword}
